% ============================================================================
% PENGESAHAN.TEX
% Halaman Pengesahan Buku Ajar
% ============================================================================

\newpage
\thispagestyle{empty}
\begin{center}
\vspace*{2cm}

\textbf{\Large HALAMAN PENGESAHAN}\\[1cm]

\textbf{BUKU AJAR}\\[0.5cm]
\textbf{PEMROGRAMAN GAME}\\[0.3cm]
\textbf{Menggunakan Unity Engine dan C\#}\\[0.3cm]
\textbf{Pendekatan Outcome-Based Education (OBE)}\\[1.5cm]

\begin{tabular}{ll}
\textbf{Nama Mata Kuliah} & : Pemrograman Game \\
\textbf{Kode Mata Kuliah} & : IF-XXX \\
\textbf{Bobot SKS} & : 3 SKS \\
\textbf{Program Studi} & : S1 Informatika \\
\textbf{Fakultas} & : Fakultas Teknologi Informasi \\
\textbf{Penulis} & : [Nama Penulis] \\
\textbf{NIP} & : [Nomor Induk Pegawai] \\
\end{tabular}\\[2cm]

Disahkan untuk digunakan sebagai buku ajar pada Program Studi S1 Informatika\\[1cm]

\vfill

\begin{tabular}{cc}
\textbf{Mengetahui,} & \textbf{Menyetujui,} \\
\textbf{Ketua Program Studi} & \textbf{Koordinator Mata Kuliah} \\[2cm]
\underline{\hspace{4cm}} & \underline{\hspace{4cm}} \\
\textbf{[Nama Ketua Prodi]} & \textbf{[Nama Koordinator]} \\
\textbf{NIP. [NIP Ketua Prodi]} & \textbf{NIP. [NIP Koordinator]} \\[1cm]
\end{tabular}

\vspace{1cm}
\textbf{[Nama Kota], [Tanggal]}

\end{center}
